\reporttitle{Prediction-Aware Coordination for Two-USV Seabed Survey Coverage under Intermittent Connectivity}
\unnumberedheading{Authors}
{\setstretch{1.1}\small
\noindent\begin{minipage}[t]{0.48\linewidth}
\centering
\textbf{Andrea Lo Iacono}\\
\texttt{andli2@kth.se}\\
\mbox{\href{https://github.com/ADreLOI}{\color{black}{\faGithubSquare}\hspace*{0.13cm}ADreLOI}}
\end{minipage}\hfill
\begin{minipage}[t]{0.48\linewidth}
\centering
\textbf{Ettore Mugisha Cirillo}\\
\texttt{emcir@kth.se}\\
\mbox{\href{https://github.com/EccirilloM}{\color{black}{\faGithubSquare}\hspace*{0.13cm}EccirilloM}}
\end{minipage}\par
\vspace{4pt}
\noindent\makebox[\linewidth][c]{\mbox{\href{https://github.com/ADreLOI/RMSW-Project}{\color{black}{\faGithubSquare}\hspace*{0.13cm}Project repository}}}\par
}
\vspace{2pt}

% Abstract describes planned work; no experimental results are claimed.
\unnumberedheading{Abstract}
This plan proposes a controlled simulation study of two unmanned surface vehicles (USVs) covering a known seabed-survey area while coordinating over an intermittent above-water wireless link. It asks when a short-horizon outage warning makes advance reservation of survey tasks useful, relative to fixed pre-allocation, reactive coordination, and a larger backlog without prediction. A simple two-dimensional coverage and burst-outage model will be developed and checked first. A persistence predictor and a signal-motion trend predictor will be compared using the same reservation action. Paired scenarios will measure mission completion, warning quality, communication load, arithmetic-operation counts, computation time, and memory use. Local coverage information will be compared with simulated coverage, while bathymetric map accuracy remains outside the model. Resource indicators will support discussion of energy-cost trends without a device-specific power model or claims about joules, battery endurance, or device lifetime. No simulator validation or experimental results are claimed at this stage.
